\cotefile{45}
\scanpage{43}
\mspage{23}
\begin{center}\uuline{BIBLIOGRAPHIE}\end{center}

[1] C.ARF. \uline{Untersuchungen über reinverzweigte Erweiterungen diskret bewerteter perfekter Körper}. Journ. Crelle, 181, 1940, p.1-44.

[2] E. ARTIN. \uline{Die gruppentheoretische Struktur der Diskriminante algebraischer Zahlkörper}. Journ. Crelle, 164, 1931, p.1-11.

[3] R. BRAUER. \uline{On the algebraic structure of group rings}. Journ. Soc. Math. Japan, 3, 1951, p.237-251.

[4] C. CHEVALLEY et A.WEIL. \uline{Über das Verhalten der Integrale} 1 \uline{Gattung bei Automorphismen des Funktionenkörpers}. Hamburg Abh., 10, 1954, p.358-361.

[5] M. DEURING. \uline{Algebren}. Ergebnisse der Math., n°IV-1, Springer, 1935.

[6] M. DEURING. \uline{Die Typen der Multiplikatorenringe elliptischer Funktionenkörper}. Hamburg Abh., 14, 1941, p.197-272.

[7] H. HASSE. \uline{Normenresttheorie galoisscher Zahlkörper mit Anwendungen auf Führer und Diskriminante abelscher Zahldörper}. Journ. Fac. Sci. Tokyo, 2, 1934, p.477-498.

[8] J.HERBRAND. \uline{Sur la théorie des groupes de décomposition, d'inertie, et de ramification}. Journ. Liouville, 10, 1931, p.481-498.

[9] S.LANG. \uline{Abelian varieties}. Interscience Tracts n°7, New-York, 1959.

[10] P.ROQUETTE. \uline{Realisierung von Darstellungen endlicher nilpotenter Gruppen}. Archiv der Math., 9, 1958, p.241-250.

[11] P.SAMUEL et O.ZARISKI. \uline{Commutative Algebra} I. Van Nostrand, 1958.

[12] O.SCHILLING. \uline{Über die Darstellungen endlicher Gruppen}. Journ. Crelle, 174, 1936, p.18.

[13] J.-P.SERRE. \uline{Corps locaux et isogénies}. Sém.BOURBAKI, 1959, n°185.

[14] R.SWAN. \uline{Induced representations and projective modules}. A paraître prochainement.

[15] A.WEIL. \uline{L'avenir des mathématiques}. Cahiers du Sud, Marseille, 1947, ainsi que Bol. \unsure{S\={a}o|São} Paulo, 1, 1946, p.55-68.

[16] A.WEIL. \uline{Sur les courbes algébriques et les variétés qui s'en déduisent}. Paris, Hermann, 1948.

[17] E.WITT. \uline{Schiefkörper über diskret bewerteten Körpern}. J. Crelle, 176, 1936, p.153-156.

[18] E.WITT. \uline{Die algebraische Struktur des Gruppenringes einer endlichen Gruppe über einem Zahlkörper}. Journ.Crelle, 190, 1952, p.231-245.

\begin{center}***\end{center}
% archive sheet number: 42

\scanpage{63}
\mspage{106}
\begin{flushright}n° 371\end{flushright}

$F=\handwritten{\prod}_{n}E_n$. Montrons qu'on peut définir sur $F$ une structure d'anneau qui induise sur $E$ la structure d'anneau donnée ; en effet, tout élément de $F$ est de la forme $(u_n)$, où $u_n\handwritten{\in}E_n$ ; on définit sur $F$ une multiplication en posant $(u_n)(v_n)=(w_n)$, avec
\[
w_n=\handwritten{\sum}_{p+q=n}u_pv_q
\tag*{(28 bis)}
\]
ce qui a un sens quel que soit $n$, puisque le second membre n'a qu'un nombre fini de termes ; l'élément $1\handwritten{\in}E_{\unsure{o|0}}$ de $E$ est évidemment élément unité pour cette loi et l'associativité se vérifie immédiatement ; on a donc ainsi sur $E$ une structure d'anneau (commutatif si $E$ est commutatif). On peut de même définir sur $B=\handwritten{\prod}_{n}A_n$ une structure d'anneau commutatif ; enfin, on définit sur $F$ une structure de $B$-\uline{algèbre} de façon analogue : si $(\handwritten{\alpha}_n)$ est un élément de $B$, on posera $(\handwritten{\alpha}_n)(u_n)=(w'_n)$ avec
\[
w'_n=\handwritten{\sum}_{p+q=n}\handwritten{\alpha}_pu_q.
\tag*{(28 ter)}
\]

Nous dirons que $B$ (resp. $F$) est la \uline{complétion canonique} de l'anneau gradué $A$ (resp. de la $A$-algèbre graduée $E$).

\begin{quote}
\handwritten{*}Cette terminologie se justifie en considérant sur $B$ (resp. $E$) la topologie \uline{produit} des topologies discrètes sur les $A_n$ (resp. $E_n$) ; alors on voit aussitôt que la multiplication dans $F$ (resp. la loi externe de $B$-module sur $F$) est le \uline{prolongement par continuité} à $F\handwritten{\times}F$ (resp. $B\handwritten{\times}F$) de la multiplication $E\handwritten{\times}E\handwritten{\to}E$ (resp. de la loi externe $A\handwritten{\times}E\handwritten{\to}E$).\handwritten{*}
\end{quote}

Soit $E'$ une seconde $A$-algèbre graduée de type $Z$ à degrés positifs, et soit $F'$ sa complétion canonique. Tout homomorphisme $f:E\handwritten{\to}E'$ de $A$-algèbres graduées se prolonge alors de façon canonique en un homomorphisme $g:F\handwritten{\to}F'$ de $B$-algèbres ; il suffit de poser, pour tout $u=(u_n)\handwritten{\in}F$
\[
g(u)=(f(u_n))
\]
et on vérifie aussitôt qu'on définit bien ainsi un homomorphisme de $B$-algèbres.
% archive sheet number: 62

\scanpage{65}
\mspage{123}
\begin{flushright}n° 371\end{flushright}

$(i_h)_{1\handwritten{\leq}h\handwritten{\leq}p}$ est la suite strictement croissante des éléments de $H$.

PROPOSITION 6.- \uline{Soit} $\unsure{K|M}$ (resp. $N$) \uline{un} $A$-\uline{module ayant une base} $(e_i)_{1\handwritten{\leq}i\handwritten{\leq}n}$ (resp. $(f_j)_{1\handwritten{\leq}j\handwritten{\leq}m}$) \uline{de} $n$ (resp. $m$) \uline{éléments} ; \uline{soit} (pour $p\handwritten{\leq}\inf(m,n)$) $(e_K)$ (resp. $(f_H)$) \uline{la base correspondante de} $\overset{p}{\handwritten{\bigwedge}}M$ (resp. $\overset{p}{\handwritten{\bigwedge}}N$). \uline{Soit} $u:M\handwritten{\to}N$ \uline{une application} $A$-\uline{linéaire}, \uline{et désignons par} $\underline{X}$ \uline{la matrice de} $u$ \uline{par rapport aux bases} $(e_i)$ \uline{et} $(f_j)$. \uline{Alors la matrice de} $\overset{p}{\handwritten{\bigwedge}}u$ \uline{par rapport aux bases} $(e_K)$ \uline{et} $(f_H)$ \uline{est la matrice} $(X_{H,K})$ \uline{à} $\binom{m}{p}$ \uline{lignes et} $\binom{n}{p}$ \uline{colonnes formée des mineurs d'ordre} $p$ \uline{de} $\underline{X}$.

En effet soit $(j_k)_{1\handwritten{\leq}k\handwritten{\leq}p}$ la suite strictement croissante des éléments de $K\handwritten{\subset}\unsure{\handwritten{[}1,n\handwritten{]}|\handwritten{[\![}1,n\handwritten{]\!]}}$ ; on a par définition
\[
\overset{p}{\handwritten{\bigwedge}}u(e_K)=u(e_{j_1})\handwritten{\wedge}u(e_{j_2})\handwritten{\wedge}\ldots\handwritten{\wedge}u(e_{j_p})\ ;
\tag*{(10)}
\]
donc l'élément de la matrice de $\overset{p}{\handwritten{\bigwedge}}u$ qui se trouve dans la ligne d'indice $H$ et la colonne d'indice $K$ est la composante d'indice $H$ du $p$-vecteur décomposable $u(e_{j_1})\handwritten{\wedge}\ldots\handwritten{\wedge}u(e_{j_p})$ ; il est par suite égal à $X_{H,K}$ en vertu de la prop.5.

On dit que la matrice $(X_{H,K})$ est la \uline{puissance extérieure} $p$-\uline{ème} de $\underline{X}$ et on la note $\overset{p}{\handwritten{\bigwedge}}\underline{X}$.

4. \uline{Développements d'un déterminant}.

Reprenons la formule (4) du n°1 donnant le déterminant d'une matrice carrée $\underline{X}=(\handwritten{\xi}_{ij})$ d'ordre $n$ ; pour toute partie $H$ de l'ensemble $\handwritten{[}1,n\handwritten{]}$, désignons par $H'$ la partie complémentaire de $H$ dans $\handwritten{[}1,n\handwritten{]}$ ; si $(i_k)_{1\handwritten{\leq}i\handwritten{\leq}p}$ (resp. $(j_h)_{1\handwritten{\leq}h\handwritten{\leq}n-p}$) est la suite strictement croissante des éléments de $H$ (resp. $H'$), on peut écrire (\S\ 3, n°2, lemme 3)
\[
x_1\handwritten{\wedge}x_2\handwritten{\wedge}\ldots\handwritten{\wedge}x_n
=\handwritten{\rho}_{H,H'}(x_{i_1}\handwritten{\wedge}x_{i_2}\handwritten{\wedge}\ldots\handwritten{\wedge}x_{i_p})\handwritten{\wedge}(x_{j_1}\handwritten{\wedge}x_{j_2}\handwritten{\wedge}\ldots\handwritten{\wedge}x_{j_{n-p}})
\tag*{(11)}
\]
où $\handwritten{\rho}_{H,H'}$ a la signification définie au \S\ 5, n°3, formule (12). Avec les notations du n°3, on a
% archive sheet number: 64

\scanpage{72}
\begin{flushright}\handwritten{22.7.961}\end{flushright}

\begin{center}Commentaires sur l'article de A.Delzant sur les formes quadratiques.\end{center}

Du point de vue rédaction;

a) Il est très confusant de faire semblant de traiter en même temps la carc. $\handwritten{\ne}2$ et la car. 2. Pour l'instant, même les structures qui entrent en jeu sont d'espèce différente, et le traitement des deux cas est vraiment différent. Si donc tu n'arrives pas à une théorie commune satisfaisante, laisse tomber la car.2, ou met\handwritten{s} dans un numéro à part ce que\deleted{\unsure{x}} tu sais en dire.

b) Contrairement à ce que suggère la lecture de l'article, le rôle des $\handwritten{\lambda}$-structures dans la théorie est pratiquement nul, dût au fait sans doute que la filtration naturelle de l'anneau des classes de formes quadratiques n'est autre que celle définie par les puissances de l'idéal d'augmentation. On récupère trivialement la $\handwritten{\lambda}$-structure sur la forme \unsure{explicite} \unsure{dudit} anneau, et il suffit de le faire à la fin.

\deleted{\unsure{a)|c)}} Du point de vue math., je suggère comme programme minimum raisonnable le suivant, qui est à portée immédiate:

1) \uline{Structure de l'anneau} $Q(A)$ \uline{des classes de formes quadratiques sur} $A$, ou pour l'instant $A$ va être un corps de caractéristique $\handwritten{\ne}2$. Pour $a\handwritten{\in}A^{\handwritten{*}}$, soit $\handwritten{\varphi}_a$ la classe de la forme quadratique $\handwritten{at^2}_{a}$, elle ne dépend que de $a\bmod A^{\handwritten{*2}}$, d'où un homomorphisme de
\[
D(A)=A^{\handwritten{*}}/A^{\handwritten{*2}}=H^{\handwritten{1}}(A,\unsure{\underline{Z}|\mathbb{Z}}/2\underline{Z})
\]
dans $Q^{\deleted{\unsure{1|2}}}(A)^{\handwritten{*}}$, notée aussi $a\handwritten{\to\varphi}_a$, d'où un homomorphisme
\[
\handwritten{\varphi:}\ \underline{Z}(D)\handwritten{\longrightarrow}Q(A)
\tag*{\handwritten{(*)}}
\]
où $\underline{Z}(D)$ est l'algèbre du groupe $D=D(A)$, à coefficients dans $\underline{Z}$. On prouvera alors le théorème suivant:

\uline{L'homomorphisme précédent est surjectif}, \uline{son noyau est} \deleted{formé du} \uline{le sous-groupe engendré par les} $e_a+e_b-e_{a'}-e_{b'}$, \uline{pour} $a,b,a',b'\handwritten{\in}D$ \uline{tels que} $\handwritten{\varphi}_a+\handwritten{\varphi}_b=\handwritten{\varphi}_{a'}+\handwritten{\varphi}_{b'}$. D'après un théorème de Witt, cette dernière relation équivaut aussi à dire que les formes des deux membres ont mêmes invariants, ce qui s'écrit plus simplement
\reversemarginpar
\margin{\handwritten{\begin{tikzpicture}[baseline=0pt]\draw[gray] (0.4,-1.7)--(0.37,0.2)--(0.4,1.6);\node at (0,0) {?};\end{tikzpicture}}}
\normalmarginpar
\[
a+b=a'+b'\ ,\quad a.b=a'.b'
\]
où $D$ est écrit additivement, et où le produit $a.b$ est le cup-produit, à valeurs dans
\[
B=H^{\handwritten{2}}(A,\underline{Z}/2\underline{Z})
\]
% archive sheet number: 71

\scanpage{73}
considéré comme sous-groupe du groupe de Brauer formé des éléments annulés par 2. Donc \uline{la structure de} $Q(A)$ \uline{est connue quand on connait les groupes} $H^{\handwritten{i}}(A,\unsure{\underline{Z}|\mathbb{Z}}/2\underline{Z})$ \uline{pour} $i=1,2$, \uline{soient} $D$ \uline{et} $B$, \uline{et le cup-produit} $D\times D\handwritten{\to}B$. L'énoncé précédent montre \deleted{en effet}\inserted{\handwritten{\unsure{même|mieux}}} que le noyau de l'homomorphisme surjectif (\handwritten{*}) est aussi le sous-groupe engendré par les éléments
\[
e_{a+u+v}-e_{a+u}-e_{a+v}+e_a
\]
pour les éléments $a,u,v$, dans $D$ tels que $uv=0$.

On en conclut en particulier que le gradué associé à $Q(A)$ est isomorphe au quotient de l'algèbre symétrique de $D$ sur $\underline{Z}$, par \deleted{\unsure{un}} l'idéal engendré par les éléments d'ordre 2 de la forme $ab-a'b'$, où $a,b,a',b'$, sont comme ci-dessus, ou ce qui revient au, même, engendré par les éléments $uv$ de $\mathrm{Symm}^{\handwritten{2}}(D)$ pour $u,v\quad D$ tels que $u.v=0$ dans $B^{\deleted{\unsure{x}}}$. (N.B. il faudrait éclaircir si $\mathrm{grad}^{\handwritten{2}}Q(A)$ \deleted{\unsure{est}} $\handwritten{\to}B$ est \deleted{\unsure{très}}\handwritten{\unsure{tjs|tjrs}} injectif, en d'autres termes si le noyau de $\mathrm{Symm}^{\handwritten{2}}(D)\quad B$ est engendré par des éléments de la forme \deleted{\unsure{uv}} $uv$).

On conclut du \deleted{\unsure{tout|notre}} théorème de structure que si ${\unsure{\mathbb{H}|H}}^{\handwritten{\bullet}}$ est un anneau gradué à degrés positifs, alors la donnée d'un homomorphisme \deleted{additif}\inserted{\handwritten{\unsure{d'anneaux augmentés}}} de $Q(A)$ dans $\handwritten{\underline{Z}\times(1+}\overset{\handwritten{\scriptscriptstyle\wedge}}{\mathbb{H}}^{\handwritten{+}}\handwritten{)}$ équivaut à la donnée d'un homomorphisme
\[
w^{\handwritten{1}}:D\handwritten{\longrightarrow}\mathbb{H}^{\handwritten{1}}
\]
tel que $W^{\handwritten{1}}(a)W^{\handwritten{1}}(b)=W^{\handwritten{1}}(a')W^{\handwritten{1}}(b')$ lorsque $a,b,a',b'$ sont comme ci-dessus, ou ce qui revient au même, tel que $W^{\handwritten{1}}(u)W^{\handwritten{1}}(v)=0$ si $u.v=0$ dans $B$; \deleted{\unsure{on}} d'ailleurs $\mathrm{grad}\ Q(A)$ joue un rôle universel pour ses propriétés, i.e. la donnée de $W$ équivaut aussi à la donnée d'un homomorphisme \deleted{\unsure{addi}} d'anneaux gradués
\[
\mathrm{grad}\ Q(A)\handwritten{\to}H^{\handwritten{\bullet}}\ .
\]
Le choix le plus naturel pour $H^{\handwritten{\bullet}}$, du point de vue ``topologique'', semble
\[
H_{\deleted{\unsure{i|1}}}^{\handwritten{\bullet}}=H^{\handwritten{\bullet}}(A,\underline{Z}/2\underline{Z})\ ,
\]
muni du cup-produit, ce qui nous donne des classes de \uline{Stiefel-Whitney}
\[
W^{\handwritten{i}}:Q(A)\handwritten{\to}H^{\handwritten{\unsure{1|i}}}(A,\underline{Z}/2\underline{Z})\ .
\]
Serre suggère de \deleted{\unsure{prendre}} considérer aussi une définition plus fine des classes de Stiefel-Whitney, en remplaçant \deleted{\unsure{le groupe|le grou}} l'extension séparable algébrique maximale de $A$ par la sous-extension réunion des sous-extensions finies de rang une puissance de 2; cela donnerait déjà un $H^{\handwritten{2}}$ ``plus fin'', i.e. s'envoyant dans le $H^{\handwritten{2}}$ habituel ...
% archive sheet number: 72

\scanpage{74}
Bien entendu, $Q(A)$ est un $\handwritten{\lambda}$-anneau, et sa $\handwritten{\lambda}$-structure se lit sur la structure d'anneau, puisqu'on aura
\[
\handwritten{\lambda}_t(\handwritten{\varphi}_a)=1+\handwritten{\varphi}_a t\ ,
\]
D'autre part, un homomorphisme d'anneaux augmentés $Q(A)\handwritten{\to}\unsure{\underline{Z}|\mathbb{Z}}\times(1+\overset{\handwritten{\scriptscriptstyle\wedge}}{A}^{\handwritten{+}})$ est automatiqueme\unsure{j|n}t un $\handwritten{\lambda}$-homomorphisme. Je ne crois pas qu'il y ait lieu de dire plus des $\handwritten{\lambda}$-structures (sauf peut-être en remarque que la filtration mise sur $Q(A)$ est celle déduite de sa $\handwritten{\lambda}$-structure).

2) \uline{Cas où} $A$ \uline{est un anneau local}, de corps résiduel $k$ tel que $\mathrm{car}\ k\handwritten{\ne}2$. Il semble bien que tout ce qui précède marqhe dans ce cas; bien entendu, $H^{\handwritten{\bullet}}(A)$ désigne la cohomologie du ``groupe fondamental'' de $\mathrm{Spec}(A)$, au sens de mon séminaire. Il faudrait alors tout rédiger dans le cadre: formes quadratiques sur un anneau local. N.B. une telle forme est dite non dégénérée si la réduite est non dégénérée; en principe, il \deleted{\unsure{faudrait être ca|faudrait être ça}} devrait être possible de façon assez triviale de relever les constructions habituelles sur $k$, en des constructions sur $A$.

Bien entendu, il serait très agréable de dépasser le programme minimum et d'arriver à une théorie des classes de Stiefel-Whitney pour des Modules quadratiques sur un \inserted{pré}schéma quelconque $S$, \deleted{\unsure{On pourrait commencer|On pourrait commen}} en supposant au besoin que les caractéristiques résiduelles sont $\handwritten{\ne}2$. \deleted{\unsure{trit|triv}} A titre de test, on peut toujours essayer de prouver l'invariance de la définition triviale relative à un Module\inserted{ quadratique} qui splitte en modules\inserted{ quadratiques} de rang 1.

Comme prime au lecteur, on peut en tous cas donner la compatibilité de la définition purement algébrique des classes de Stiefel-Whitney, et de la définition topologique, lorsque $A$ est l'anneau local\inserted{ d'un point $s$} d'un schéma de type fini\inserted{$S$} sur $\underline{C}$. Un module quadratique sur $A$ définit un module quadratique sur un voisinage de $S$, donc $S$ si on veut, i.e. un fibré principal homogène algébrique sous le schéma en groupes orthogonal. Par Gaga il lui correspond un fibré principal analytique, donc topologique, d'où des classes de Stiefel-Whitney topologiques dans $H^{\handwritten{\bullet}}(S_{\mathrm{top}},\underline{Z}/2\underline{Z})$. Or on a un homomorphisme canonique,que tu définiras:
\[
H^{\handwritten{\bullet}}(S_{\mathrm{top}},\underline{Z}/2\underline{Z})\handwritten{\longrightarrow}H^{\handwritten{\bullet}}(A,\underline{Z}/2\unsure{\underline{\mathbb{Z}}|\underline{Z}})\ ,
\]
et on voit aussitôt que les images des classes caractéristiques topologiques sont les classes caractéristiques algébriques, (car il suffit de le vérifier en rang 1).
% archive sheet number: 73

\scanpage{75}
3) \uline{Essayer de trouver des définitions raisonnables en car.}2, \uline{ou de préférence \unsure{indépendantes} de la caractéristique}. Une théorie vraiment satisfaisante devrait marcher sur tout préschéma $S$, et réciproquement. Je suis loin d'avoir des idées claires sur la question, et me borne à des remarques vagues:

a) Lorsqu'on est sur un corps de caractéristique \deleted{2}\inserted{quelconque}, on est sans doute obligé de tenir compte des formes quadratiques dégénérées, et de définir un \unsure{ranneau|anneau} $Q'(k)$ comme défini par les classes de formes quadratiques, dégénérées ou non. \handwritten{[}Dans le cas d'un corps de \deleted{\illegible{1}} caractéristique $\handwritten{\ne}2$, ce serait $Q(k)+\unsure{\underline{Z}|\mathbb{Z}}$, où le \deleted{deuxième facteur $\underline{Z}$ est un} multiplication est donnée par $(\unsure{g|f},m).(g,n)=(fg,\deg(f)n+\deg(g)m+mn)$, le facteur en rabiot $\underline{Z}$ correspondant aux formes nulles, sur des vectoriels de rang quelconque\handwritten{]}.

b) En caractéristique quelconque, garder à l'esprit la distinction entre le shhéma en groupes $(\underline{Z}/2\underline{Z})_S$ et le shhéma en groupes $(\handwritten{\mu_2})_S$ sur $S$, le dernier étant le noyau de la puissance deuxième dans le groupe multiplicatif $\underline{G}_{m\,S}$. Si $\unsure{S|A}$ est local d'anneau $A$, alors la suite exacte de cohomologie donne $H^{\handwritten{1}}(S,\handwritten{\mu_2})=A^{\handwritten{*}}/A^{\handwritten{*2}}$, $H^{\handwritten{2}}(S,\handwritten{\mu_2})={}_{\handwritten{2}}H^{\handwritten{2}}(A,\underline{G}_m)$, éléments annulés par 2 dans le groupe de Brauer, par contre $H^{\handwritten{1}}(S,\underline{Z}/2\underline{Z})=A/\handwritten{\wp}A$ lorsque $2A=0$. Il faut donc s'attendre à trouver des classes de Stiefel-Whitney de type mixte, faisant à la fois intervenir les deux groupes voisins de coefficients. Cette remarque se précise quand on regarde les schémas en groupes orthogonal\handwritten{s}, sur $\underline{Z}$ disons: en dimension paire, c'est l'extensions de $\underline{Z}/2\underline{Z}$ par un schéma en groupes à fibres connexes (le schéma spécial-orthogonal), par exemple en dim 2 c'est \deleted{\unsure{l'extension}} un produit semi-direct de $\underline{Z}/2\underline{Z}$ par $\underline{G}_m$ sur lequel il opère par symétrie; en dimension impaire, c'est une extension de $\handwritten{\mu_2}$ par le groupe spécial-orthogonal, par exemple en dimension 1, c'est $\handwritten{\mu_2}$. D'où suivant la parité, comme premier invariant ``discret'' d'un fibré principal \unsure{orthhogonal|orthogonal}, une 1-classe de groupe $\underline{Z}/2\underline{Z}$,ou $\handwritten{\mu_2}$.
% archive sheet number: 74

\scanpage{80}
\mspage{105}
\begin{flushright}n° 371\end{flushright}

Le cas le plus important est celui où $M$ est un $A$-module \uline{libre de dimension finie} $n$ ; on identifie alors canoniquement $\handwritten{\bigwedge}(M^{\handwritten{*}})$ et $(\handwritten{\bigwedge}(M))^{\handwritten{*}\mathrm{gr}}$ au moyen de l'isomorphisme canonique $\handwritten{\lambda}_M$ ; un multivecteur (resp. $p$-vecteur) sur $M^{\handwritten{*}}$ est donc identifié à une forme linéaire sur $\handwritten{\bigwedge}(M)$ (resp. $\overset{p}{\handwritten{\bigwedge}}M$) et appelé une \uline{multiforme} (resp. une $p$-\uline{forme}). \handwritten{*}Avec cette identification, la formule (22) s'écrit
\[
x_1\handwritten{\wedge}x_2\handwritten{\wedge}\ldots\handwritten{\wedge}x_n\ ,\ x_1^{\handwritten{*}}\handwritten{\wedge}x_2^{\handwritten{*}}\handwritten{\wedge}\ldots\handwritten{\wedge}x_n^{\handwritten{*}}=\det(\handwritten{\langle}x_{\unsure{1|i}}\ ,\ x_j^{\handwritten{*}}\handwritten{\rangle}).\ \handwritten{*}
\tag*{(27)}
\]

Si $(e_{\unsure{i|1}})_{1\handwritten{\leq}i\handwritten{\leq}n}$ est une base de $M$, $(e_i^{\handwritten{*}})$ la base duale dans $M^{\handwritten{*}}$, les éléments
\[
e_H^{\handwritten{*}}=e_{i_1}^{\handwritten{*}}\handwritten{\wedge}e_{i_2}^{\handwritten{*}}\handwritten{\wedge}\ldots\handwritten{\wedge}e_{i_p}^{\handwritten{*}}
\tag*{(28)}
\]
(avec les notations de la prop.3) forment une base de $\handwritten{\bigwedge}(M^{\handwritten{*}})$ lorsque $H$ parcourt les parties de $\unsure{\handwritten{[}1,n\handwritten{]}|\handwritten{[\![}1,n\handwritten{]\!]}}$ ; cette base n'est autre que la \uline{base duale} de la base $(e_H)$ de $\handwritten{\bigwedge}(M)$, car il résulte de (27) que si $H$, $K$ sont deux parties distinctes de $p$ éléments de $\handwritten{[}1,n\handwritten{]}$, $(i_h)_{1\handwritten{\leq}h\handwritten{\leq}p}$, $(j_h)_{1\handwritten{\leq}h\handwritten{\leq}p}$ les suites strictement croissantes des éléments de $H$ et de $K$ respectivement, on a
\[
\handwritten{\langle}e_{i_1}\handwritten{\wedge}e_{i_2}\handwritten{\wedge}\ldots\handwritten{\wedge}e_{i_p}\ ,\ e_{j_1}^{\handwritten{*}}\handwritten{\wedge}e_{j_2}^{\handwritten{*}}\handwritten{\wedge}\ldots\handwritten{\wedge}e_{j_p}^{\handwritten{*}}\handwritten{\rangle}=0
\]
puisqu'il existe un indice $k$ \unsure{tel|til} que $j_k\handwritten{\notin}H$, donc $\handwritten{\langle}e_{i_h}\ ,\ e_{j_k}^{\handwritten{*}}\handwritten{\rangle}=0$ quel que soit $h$.

Si $M$ et $N$ sont deux $A$-modules libres de type fini, $u:M\handwritten{\to}N$ une application $A$-linéaire, ${}^{t}(\handwritten{\bigwedge}(u))$ s'identifie canoniquement à $\handwritten{\bigwedge}({}^{t}u)$ en vertu de la commutativité du diagramme (25) ; il en résulte que si $u$ est \uline{bijectif}, le contragrédient ${}^{t}(\handwritten{\bigwedge}(u))^{-1}$ s'identifie à $\handwritten{\bigwedge}({}^{t}u^{-1})$.

\uline{Remarque}.- A) \uline{Rabiot qui doit remonter au} \S3 (\uline{Algèbres graduées}), \uline{où il a été omis par inadvertance}.

Soient $A$ un anneau commutatif gradué de type $Z$ à degrés positifs, $E$ une $A$-algèbre graduée de type $Z$ à degrés positifs, $(A_n)$ et $(E_n)$ les graduations respectives de $A$ et de $E$. En tant que groupe additif, $E$ est un sous-groupe du \uline{produit}
% archive sheet number: 79 (preceded by a crossed-out unreadable archive number)
